\documentclass[11pt,a4paper]{article}
\usepackage{turkos}

\begin{document}
\phaseheader{9}{Synchronization and Concurrency in the Kernel}{Now that tasks can be interrupted anywhere, make shared data safe: interrupt control, spinlocks, wait queues, mutexes and semaphores}{56}{62}

\ataglance{2 weeks}{4}{Phase 8: preemptive kernel threads, \code{schedule()}, sleep and wake-up}{Interrupt-safe critical
sections, spinlocks built on the atomic \code{xchg} instruction, wait queues with lost-wakeup-free sleeping, mutexes,
counting semaphores and condition variables, a keyboard driver that blocks instead of polling, the classic
concurrency problems as kernel tests, and a written lock-ordering rule.}

\section{Why this phase}

Phase 8 introduced a new kind of bug. A task can now be interrupted between any two instructions, and another task
or an interrupt handler can run and touch the same data. Before preemption, the kernel's lists and counters were
only changed by one flow of control at a time. Now the run queue, the heap's free list, the PMM bitmap, the console
cursor and the keyboard buffer can all be modified halfway through an update, and the result is corruption that
appears rarely, at random, and never under the debugger.

This phase gives you the tools to prevent that, and the discipline to use them everywhere. It follows \OSTEP's
second big idea, concurrency, closely: build locks from a hardware atomic instruction, then build sleeping primitives
on top of the locks, then use them to solve the classic problems. You will also think about deadlock for the first
time, because once code takes more than one lock, the order it takes them in starts to matter.

\section{Concepts you need}

\subsection{Race conditions and critical sections}

A \textbf{race condition} is a bug whose outcome depends on the timing of concurrent events. \OSTEP chapter 26 uses
the simplest possible example: two threads running \code{counter = counter + 1}. The statement compiles to three
instructions (load, add, store), and a switch between the load and the store loses an update (Table~\ref{tab:race}).
The code that touches shared data is a \textbf{critical section}, and the goal is \textbf{mutual exclusion}: at most
one flow of control inside it at a time.

\begin{table}[htbp]
\centering\small\renewcommand{\arraystretch}{1.25}
\begin{tabular}{@{}c l l c c c@{}}
\toprule
\tkth{Step} & \tkth{Task A} & \tkth{Task B} & \tkth{A's EAX} & \tkth{B's EAX} & \tkth{counter} \\
\midrule
1 & \code{mov eax, [counter]} & & 50 & -- & 50 \\
2 & \code{add eax, 1} & & 51 & -- & 50 \\
3 & \multicolumn{2}{l}{\textit{timer interrupt: switch to B}} & 51 & -- & 50 \\
4 & & \code{mov eax, [counter]} & 51 & 50 & 50 \\
5 & & \code{add eax, 1} & 51 & 51 & 50 \\
6 & & \code{mov [counter], eax} & 51 & 51 & 51 \\
7 & \multicolumn{2}{l}{\textit{switch back to A}} & 51 & 51 & 51 \\
8 & \code{mov [counter], eax} & & 51 & 51 & \textbf{51} \\
\bottomrule
\end{tabular}
\caption{Two increments, one lost. Each task has its own saved EAX, so the switch at step 3 is invisible to A.}
\label{tab:race}
\end{table}

\subsection{Where concurrency comes from in Turk-OS}

On a single CPU there are exactly two sources of interleaving: \textbf{interrupts}, which can run a handler between
any two instructions of a task, and \textbf{preemption}, which the timer interrupt causes. So on one CPU,
\textbf{disabling interrupts} around a critical section is a complete lock (\OSTEP\,\S28.5). It is cheap and simple,
but it doesn't scale: with several CPUs, the other CPUs keep running. That is why Turk-OS wraps interrupt control in
a real \textbf{spinlock} built on an atomic instruction from the start. On one CPU the spinning never happens, but the
code is already correct for the SMP track in Phase 15.

\subsection{Atomic instructions and spinlocks}

A lock needs an operation that reads and writes memory in one indivisible step. On x86, \code{xchg} with a memory
operand is always atomic: it swaps a register with memory while locking out other CPUs. That makes it a
\textbf{test-and-set}: write 1 and get the old value back; if the old value was 0, you have the lock
(\OSTEP\,\S28.7, \OSC\,\S6.4). GCC exposes it as \code{__atomic_exchange_n}. Other atomics you will meet are
compare-and-swap (\code{lock cmpxchg}) and fetch-and-add (\code{lock xadd}).

\subsection{Sleeping instead of spinning}

Spinning wastes the CPU, and on one CPU it can never succeed while the holder isn't running. When a task has to wait
for something that may take a while (a key press, a disk block, a free slot in a buffer), it should \textbf{block}:
go on a \textbf{wait queue}, let the scheduler run someone else, and be put back on the ready queue by whoever
makes the condition true. Two rules make this correct. First, going to sleep must be atomic with respect to the
wake-up, or you get the \textbf{lost wake-up} problem (\MOS\,\S2.3.4): the task checks the condition, the wake-up
happens, and then the task sleeps forever. Second, a woken task must re-check its condition in a \code{while} loop,
because by the time it runs, someone else may have taken what it was waiting for (\OSTEP chapter 30 calls this Mesa
semantics).

\begin{conceptbox}{Three primitives, three jobs}
\textbf{Spinlock}: short critical sections, and the only lock allowed in interrupt handlers. The holder must never
sleep. \textbf{Mutex}: a sleeping lock with an owner, for longer critical sections in task context, such as the
file system in Phase 13. Never from an interrupt handler. \textbf{Semaphore}: a counter with a wait queue
(\OSTEP chapter 31). \code{down} waits until the count is positive and decrements it; \code{up} increments it and
wakes one waiter. It serves for counting resources and for signalling from an interrupt handler (which may call
\code{up} but never \code{down}). \textbf{Condition variables} (\code{wait}, \code{signal}, \code{broadcast}) with a
mutex round out the set.
\end{conceptbox}

\subsection{Deadlock}

A \textbf{deadlock} is a set of tasks each waiting for something another one holds. \OSTEP chapter 32, \MOS chapter 6
and \OSC chapter 8 list the four conditions that must all hold for it: mutual exclusion, hold-and-wait, no preemption,
and circular wait. Breaking any one prevents deadlock, and in kernels the practical choice is to break circular wait
with a global \textbf{lock order}: if every piece of code takes locks in the same order, no cycle can form. \MOS\,\S6.3
also names the ``ostrich algorithm'', ignoring the problem, and explains when that is a reasonable choice (for the
kernel, it isn't).

\section{Reading plan}

\begin{readingplan}
\must & \OSTEP & Ch.~26 \emph{Concurrency: An Introduction} & 273--288 & Races, critical sections, atomicity. \\
\must & \OSTEP & Ch.~28 \emph{Locks} (interrupts, test-and-set, compare-and-swap, spinning versus sleeping) & 301--324 & The design behind Tasks 9.1--9.2. \\
\must & \OSTEP & Ch.~30 \emph{Condition Variables}; Ch.~31 \emph{Semaphores} & 339--372 & Producer/consumer, the \code{while} rule, semaphores as locks and as signals. \\
\must & \OSTEP & Ch.~32 \emph{Common Concurrency Problems} & 373--388 & Atomicity and order violations; deadlock and how to prevent it. \\
\should & \OSC & Ch.~6 \emph{Synchronization Tools} (critical sections, Peterson, hardware support, mutexes, semaphores, monitors, liveness) & 257--288 & A second, more formal treatment. \\
\should & \OSC & \S7.1--7.2 \emph{Classic Problems}; \emph{Synchronization within the Kernel} & 289--299 & What Linux and Windows use inside their kernels. \\
\should & \MOS & \S2.3 \emph{Interprocess Communication}; \S2.5 \emph{Classical IPC Problems} & 119--148, 167--172 & Sleep and wakeup, semaphores, mutexes, monitors, message passing; philosophers. \\
\should & \MOS & \S6.1--6.6 \emph{Deadlocks} & 436--458 & Resources, conditions, detection, avoidance, prevention. \\
\should & \OSDI & \S2.2.8 \emph{Message Passing}; \S2.6.9 \emph{Interprocess Communication in MINIX 3} & 85--88, 178--182 & The alternative MINIX chose instead of kernel locks. \\
\could & \OSTEP & Ch.~27 \emph{Thread API}; Ch.~29 \emph{Lock-based Concurrent Data Structures} & 289--300, 325--338 & Locking real data structures well. \\
\could & \OSDI & \S3.3 \emph{Deadlocks} & 237--252 & The same theory with the I/O perspective. \\
\end{readingplan}

\begin{minixbox}
The MINIX 3 kernel needs almost no locks: it is small, it is not preemptible, and it runs with interrupts disabled.
Everything else (drivers, the file system, the process manager) runs as separate processes that cooperate through
\textbf{synchronous message passing}: \code{send} blocks until the receiver calls \code{receive}, so every message is
a rendezvous with no buffering (\OSDI\,\S2.2.8 and \S2.6.9). Read \code{mini_send} and \code{mini_receive} in
\texttt{kernel/proc.c (07400)}. Notice the check that refuses a \code{send} which would close a cycle of processes
sending to each other: deadlock detection built into the IPC primitive itself.
\end{minixbox}

\section{Build it}

\task{Interrupt control that nests}
Replace every bare \code{cli}/\code{sti} with a pair that saves and restores the previous state. A function that
disables interrupts must not enable them on exit if they were already disabled when it was called.

\begin{lst}{c}{include/arch/irq.h}
static inline uint32_t irq_save(void)
{
    uint32_t flags;
    __asm__ volatile ("pushf; pop %0; cli" : "=r"(flags) : : "memory");
    return flags;
}

static inline void irq_restore(uint32_t flags)
{
    if (flags & 0x200)                      /* IF was set: enable interrupts again */
        __asm__ volatile ("sti" : : : "memory");
}
\end{lst}

\task{Spinlocks}
A spinlock disables interrupts (on one CPU, that is what actually excludes everyone else) and then takes an atomic
flag, so the same code stays correct on several CPUs. Record the owner: taking a lock you already hold is a guaranteed
deadlock, and an assertion turns it into an immediate, readable panic. The \code{"memory"} clobbers and the
acquire/release orderings stop the compiler from moving accesses out of the critical section.

\begin{lst}{c}{kernel/sync/spinlock.c}
typedef struct {
    volatile uint32_t locked;
    uint32_t          saved_flags;          /* interrupt state of the holder before locking */
    struct task      *owner;                /* debugging: who holds it                      */
    const char       *name;
} spinlock_t;

void spin_lock(spinlock_t *l)
{
    uint32_t flags = irq_save();
    ASSERT(!(l->locked && l->owner == current));          /* already mine: deadlock */
    while (__atomic_exchange_n(&l->locked, 1, __ATOMIC_ACQUIRE))
        __asm__ volatile ("pause");                       /* only ever spins on SMP */
    l->saved_flags = flags;
    l->owner = current;
}

void spin_unlock(spinlock_t *l)
{
    ASSERT(l->locked && l->owner == current);
    uint32_t flags = l->saved_flags;
    l->owner = NULL;
    __atomic_store_n(&l->locked, 0, __ATOMIC_RELEASE);
    irq_restore(flags);
}
\end{lst}
\verify{Taking the same lock twice panics with the lock's name. Releasing locks in the reverse order they were
taken restores the original interrupt state at every level.}

\task{Wait queues without lost wake-ups}
A wait queue is a list of blocked tasks. \code{sleep_on(wq, lock)} is called with a spinlock held, which on one CPU
means interrupts are off. It marks the task blocked and queues it \emph{before} releasing the lock, so a wake-up can't
slip in between. Then it schedules, and it takes the lock again before returning, so the caller can re-check its
condition safely. \code{wake_up_one} and \code{wake_up_all} move tasks from the wait queue to the ready queue; they
never block, so interrupt handlers may call them.

\begin{lst}{c}{kernel/sync/waitq.c}
struct wait_queue { struct list_node tasks; };

/* Caller holds lk. Returns with lk held again, after a wake_up. */
void sleep_on(struct wait_queue *wq, spinlock_t *lk)
{
    uint32_t caller_flags = lk->saved_flags;
    current->state = TASK_BLOCKED;
    list_add_tail(&wq->tasks, &current->node);
    lk->owner = NULL;
    __atomic_store_n(&lk->locked, 0, __ATOMIC_RELEASE);  /* release, keep interrupts off */
    schedule();                                           /* returns after we are woken  */
    spin_lock(lk);
    lk->saved_flags = caller_flags;                       /* keep the caller's IF state  */
}

void wake_up_one(struct wait_queue *wq)
{
    if (!list_empty(&wq->tasks)) {
        struct task *t = container_of(wq->tasks.next, struct task, node);
        list_del(&t->node);
        make_ready(t);
    }
}
\end{lst}

\task{Semaphores, mutexes and condition variables}
Build the three sleeping primitives on spinlocks and wait queues. The semaphore is shown below. A mutex is a
semaphore with an initial count of one plus an owner field and an assertion that only the owner unlocks. A condition
variable is a wait queue used together with a mutex: \code{cv_wait} releases the mutex and sleeps atomically, and
re-acquires the mutex before returning.

\begin{lst}{c}{kernel/sync/semaphore.c}
struct semaphore {
    int               count;
    spinlock_t        lock;
    struct wait_queue waiters;
};

void sem_down(struct semaphore *s)          /* task context only: may sleep */
{
    spin_lock(&s->lock);
    while (s->count == 0)                   /* while, not if: re-check after waking */
        sleep_on(&s->waiters, &s->lock);
    s->count--;
    spin_unlock(&s->lock);
}

void sem_up(struct semaphore *s)            /* safe from interrupt handlers */
{
    spin_lock(&s->lock);
    s->count++;
    wake_up_one(&s->waiters);
    spin_unlock(&s->lock);
}
\end{lst}

\task{A keyboard that blocks}
Rewrite the keyboard buffer as a producer/consumer pair. The interrupt handler is the producer: it puts the character
in the ring buffer and calls \code{sem_up} on a semaphore that counts available characters. \code{kbd_getchar} is the
consumer: \code{sem_down}, then take a character. The \code{sti; hlt} loop from Phase 4 disappears; a task waiting for
a key now simply sleeps and costs nothing.
\verify{With the monitor waiting for input and no other task ready, \code{ps} shows the monitor as blocked and the
idle task running.}

\task{Make the classic problems into tests}
Add these to \code{make test}. Each one should run long enough (a few seconds) to expose timing bugs.

\begin{center}
\small\renewcommand{\arraystretch}{1.3}
\begin{tabularx}{\linewidth}{@{}l X@{}}
\toprule
\tkth{Test} & \tkth{What it shows} \\
\midrule
\code{race_counter} & four tasks increment a shared counter 100\,000 times each. Without a lock the total comes out short; with a spinlock it is exactly 400\,000. To make the race visible on one CPU, widen the window between load and store with a short busy loop. \\
\code{bounded_buffer} & three producers and three consumers share a 16-slot buffer, using semaphores for empty slots, full slots and mutual exclusion (\OSTEP chapter 31). Every item produced is consumed exactly once. \\
\code{philosophers} & five tasks, five forks. Taking forks in a fixed global order (lowest number first) runs without deadlock; for the lesson, run a version without ordering once and watch it freeze. \\
\code{cv_ping_pong} & two tasks alternate strictly using a mutex and a condition variable, 10\,000 rounds. \\
\bottomrule
\end{tabularx}
\end{center}
\verify{All four tests pass repeatedly. The unlocked version of \code{race_counter} fails, which proves the test
can see races.}

\task{Audit the kernel and write the lock order}
Go through every shared structure you have built and give it an owner. Write the result into a comment block in
\code{include/locks.h} and follow it from now on. Locks are always taken from top to bottom in the table, never upward.

\begin{center}
\small\renewcommand{\arraystretch}{1.25}
\begin{tabularx}{\linewidth}{@{}l l X@{}}
\toprule
\tkth{Data} & \tkth{Protected by} & \tkth{Notes} \\
\midrule
task list, run queue, sleep queue & \code{sched_lock} (spinlock) & taken inside \code{schedule()}, \code{make_ready()} and the timer handler \\
keyboard buffer & its semaphore and a spinlock & producer is the IRQ handler \\
console and \code{kprintf} & \code{console_lock} (spinlock) & also used from interrupt context, so it must be a spinlock \\
kernel heap & \code{heap_lock} (spinlock) & leaf lock: never call anything that takes another lock while holding it \\
PMM bitmap & \code{pmm_lock} (spinlock) & leaf lock \\
\bottomrule
\end{tabularx}
\end{center}

\section{Testing and debugging}

Concurrency bugs hide from debuggers because stopping the machine changes the timing. Rely on assertions instead:
check lock ownership, check that interrupts are disabled where you expect them to be (\code{ASSERT(!(read_eflags() &
0x200))}), and run tests many times. Changing the timer frequency (10\,Hz, 1000\,Hz) shuffles the interleavings and
often exposes a race that was hiding.

\begin{pitfalls}
System freezes the first time an IRQ handler runs & The handler took a mutex or called \code{sem_down}, which sleeps in interrupt context & Handlers may only take spinlocks and call \code{wake_up}/\code{sem_up}. Assert that you are not in an interrupt in every sleeping function. \\
Deadlock when \code{kprintf} is called from an interrupt & The interrupted task held \code{console_lock}, and the lock didn't disable interrupts & Spinlocks must disable interrupts: that is exactly what \code{spin_lock} does. \\
A task sleeps forever although the event happened & Lost wake-up: the condition was checked, then the wake-up came, then the task slept & Queue the task before releasing the lock, as \code{sleep_on} does. \\
Consumers take items that aren't there & \code{if} instead of \code{while} around the wait & Always re-check the condition in a \code{while} loop. \\
Interrupts stay off after some code path & An early \code{return} skipped \code{spin_unlock} or \code{irq_restore} & One exit path per locked function; assert IF at the top of the idle loop. \\
Occasional hang under load & Two locks taken in opposite orders & Follow the lock order table; see Going further for automatic checking. \\
\end{pitfalls}

\section{Milestone}

\begin{milestonebox}[Safe concurrency]
\begin{checklist}
  \item \code{irq_save}/\code{irq_restore}, spinlocks, wait queues, semaphores, mutexes and condition variables are implemented.
  \item The unlocked counter test loses updates; the locked one is exact.
  \item Bounded buffer, philosophers and ping-pong tests pass repeatedly in \code{make test}.
  \item The keyboard driver blocks the reader on a semaphore; no busy waiting remains in the kernel.
  \item Every shared structure has a documented lock, and the lock order is written down.
\end{checklist}
\end{milestonebox}

\begin{questionbox}
\begin{enumerate}
  \item Why is disabling interrupts a complete lock on one CPU but not on several?
  \item Write test-and-set with \code{xchg}. Why must the exchange be atomic, and what does ``atomic'' mean on a multiprocessor?
  \item What is the lost wake-up problem, and which line of \code{sleep_on} prevents it?
  \item Compare a spinlock, a mutex, a semaphore and a condition variable. Which may be used in an interrupt handler?
  \item Why must a woken task re-check its condition in a \code{while} loop?
  \item Name the four conditions for deadlock. Which one does a lock order break, and how?
  \item MINIX uses synchronous message passing instead of shared-memory locks. What does it gain and what does it lose?
\end{enumerate}
\end{questionbox}

\section{Going further}

Build a small \emph{lock validator}: give each lock a class number, record the order in which each task acquires
classes, and panic the first time an order inversion appears, even if it didn't deadlock that time. Linux's lockdep
works this way. Implement reader-writer locks (\OSTEP chapter 31) for data that is read far more often than written.
Or add MINIX-style synchronous messages (\code{send}, \code{receive}, \code{sendrec}) as a second IPC mechanism
between kernel threads, and compare the code you write with each style.

\nextphase{10}{User Mode and System Calls}
\booklegend
\end{document}
